\chapter*{Glossar}
\label{Glossar}

\begin{description}


\item[Fehlertolerante Programme:] Programme, die nach Systemabst"urzen oder sonstigen Problemen dort weitermachen k"onnen, wo sie aufgeh"ort haben. 

In Verbindung mit einem verteilten System meint dies meist die Toleranz des Gesamtsystems gegen"uber einzelnen Rechnerabst"urzen.



\item[Parsen:]	{\em Parsen} ist das Herausfiltern von (syntaktischen) Informationen aus einem Datenstrom. 
				
Beim Parsen von HTML-Dokumenten werden alle Hypertext-Informationen herausgefiltert. Nachdem ein HTML-Dokument geparsed wurde, stehen alle S"atze des Dokuments in einer Datenstruktur zur Verf"ugung.

Ein Beispiel erhellt den Vorgang: Die Einstiegsseite der Fachhochschule, wie sie in Abbildung \ref{FachhochschuleHTML} auf Seite \pageref{FachhochschuleHTML} zu sehen ist, wird zu den folgenden S"atzen geparsed:
\nl
{\tt \nl
URL: http://www.informatik.fh-muenchen.de/welcome.html	\nl
SEN: Dies ist der Hypermedia-Informationsserver des Fachbereiches 07 an der Fachhochschule M"unchen	\nl
SEN: Das Informationssystem ist ein Teil des internationalen World Wide Web	\nl
}
				
Die "Uberschrift des HTML-Dokuments wurde nicht als Satz erkannt, da diese das gestelltes Kriterium: ein Satz mu"s mindestens acht W"orter enthalten, nicht erf"ullt.
				
\begin{figure}[p]
	\makebox[\textwidth]{\epsfxsize=\textwidth \epsffile{FachhochschuleHTML.eps} }
	\caption{Die Webseite der Fachhochschule M"unchen in einem Webbrowser dargestellt.}
	\label{FachhochschuleHTML}
\end{figure}


\item[Proxy-Server:]	Proxy-Server oder auch Web-Caches speichern alle von Benutzern des Internet geholten Hypertextdokumente. Damit werden Seiten, die von mehrmals angefragt werden, nur aus dem Speicher des "`nahen"' Proxies geholt und m"ussen nicht erst langwierig vom Internet geladen werden.


\item[Regul"are Ausdr"ucke:]	Mit regul"aren Ausdr"ucken kann man in algebraischer Form, Suchmuster definieren.[Foundations]

Einfach ausgedr"uckt kann man mit reg"ularen Ausdr"ucken Suchmuster sehr leicht definieren. So benutzen viele Unix-Werkzeuge regul"are Ausdr"ucke um Dateinamen allgemein zu bennenen. In der Programmiersprache Perl bekommt man regul"are Ausdr"ucke gleich im Sprachumfang mitgeliefert, weshalb sich diese Sprache auch gut f"ur Textparsing eignet.

Will man z.B. alle auseinandergezogenen Bindestrichw"orter zu Mehrworten ohne Bindestrich konvertieren, dann kann man dies durch den Ausdruck:
\nl 
{\tt s/$\backslash$s-/ /g;	\nl
}
erledigen. Im Klartext: Ersetze (s) jedes Zeichenpaar, welches aus einem Leerzeichen oder Tabulator besteht ($\backslash$s) und auf den ein Bindestrich folgt (-) durch ein Leerzeichen ( ). F"uhre diese Ersetztung f"ur jedes Zeichen in der Zeichenkette aus (g).


\item[Thrashen:] {\em engl.} dreschen, {\em hier:} um sich schlagen. 

So wird das Systemverhalten von Betriebssystemen genannt, welches auftritt, wenn Prozesse so gro"s werden, da"s andere Prozesse nicht nur ausgelagert (paging), sondern "uberdies noch selber aufgrund von Seitenverdr"angnungsalgorithmen auf die Platte ausgelagert werden. Das System arbeitet dann sozusagen die Programmbefehle nicht mehr aus dem Hauptspeicher heraus ab, sondern mu"s erst jeden Befehl von der Festplatte laden. Diese, aufgrund der Mechanik der Festplatte bis zu eine Million mal langsamere Abarbeitung der Befehle macht einen Rechner dann unbenutzbar.\cite{Denning1968}


\item[Verteilte Programme:] Programme die auf mehreren Rechnern gleichzeitig Rechenkapazit"at nutzen und dies auch koordinieren.



\item[Web-Site:] Web-Sites sind Rechner im Internet, die Hypertextdokumente zum weltweiten Abruf zur Verf"ugung stellen.
\end{description}
